% ============================================================
%  پیوست الف: سنت اسلامی–ایرانی و مفهوم آزادی
% ============================================================
\appendix
\section{پیوست الف: آزادی در سنت اسلامی–ایرانی}
\label{app:islamic}

\begin{tcolorbox}[wavebox, title={درباره‌ی این پیوست}]
این پیوست به سه بحث اصلی می‌پردازد: (۱) تقابل مفهوم آزادی اسلامی با بردگی یونانی؛ (۲) تنش آزادی و استبداد در حکمرانی اسلامی–ایرانی؛ و (۳) بحث ایرانشهری و دیدگاه‌های طبقه‌بندی‌شده درباره‌ی اسلام و آزادی.
\end{tcolorbox}

% ─────────────────────────────────────────────────────────────
\subsection{آ.۱ — آزادی در برابر بردگی: تقابل اسلام و یونان باستان}
\label{subsec:islam_slavery}
% ─────────────────────────────────────────────────────────────

\subsubsection{وضعیت بردگی در پولیس یونانی}

همان‌طور که بیان شد، آزادی یونانی (ἐλευθερία) بر پایه‌ی نظام برده‌داری ممکن بود. \thinker{ارسطو} بردگی را «طبیعی» می‌دانست: برخی انسان‌ها به‌طور طبیعی برده‌اند و قادر به خودگردانی نیستند.\footnote{\lr{Aristotle, \textit{Politics}, I, 2, 1252a--1255b.}} این نگاه آزادی سیاسی را برای اقلیت محفوظ می‌کرد و اکثریت را از آن محروم. % FIXED: \lr in footnote

\subsubsection{پاسخ اسلامی: توحید و کرامت عمومی}

اسلام از همان آغاز با رویکردی متفاوت وارد این مسئله شد. مفهوم توحیدی خداوند در اسلام پیامدی مستقیم برای آزادی داشت: اگر تنها خداوند ربّ است، پس هیچ انسانی به‌طور ذاتی «ربّ» انسان دیگری نیست.

\begin{tcolorbox}[quotebox]
«یا أَیُّهَا النَّاسُ إِنَّا خَلَقْنَاکُمْ مِنْ ذَکَرٍ وَأُنْثَى وَجَعَلْنَاکُمْ شُعُوباً وَقَبَائِلَ لِتَعَارَفُوا إِنَّ أَکْرَمَکُمْ عِنْدَ اللَّهِ أَتْقَاکُمْ» (حجرات: ۱۳)
\end{tcolorbox}

این آیه کرامت عمومی بشر را اعلام می‌کند. اگرچه اسلام بردگی را به‌طور فوری لغو نکرد، اما دو تحول مهم ایجاد کرد: (۱) \textbf{تشویق به آزادسازی برده} (\lr{itq}) به‌عنوان عبادت؛ (۲) \textbf{محدودیت منابع بردگی} به اسرای جنگ—که خود بحث‌برانگیز بود.

\subsubsection{مقایسه‌ی تطبیقی: اسلام در برابر یونان}

\begin{table}[h!]
\centering
\caption{مقایسه‌ی رویکرد اسلامی و یونانی به بردگی و آزادی}
\label{tab:islamgreece}
\renewcommand{\arraystretch}{1.6}

\begin{tabularx}{\textwidth}{|>{\bfseries}p{3cm}|>{\raggedright}X|>{\raggedright\arraybackslash}X|}
\hline
\rowcolor{PrimaryDeep}
\textcolor{white}{\textbf{معیار}} &
\textcolor{white}{\textbf{یونان باستان}} &
\textcolor{white}{\textbf{سنت اسلامی}} \\
\hline

\rowcolor{EraAncient!15}
پایه‌ی کرامت & طبیعت (فوسیس) — طبیعی‌بودن بردگی & توحید — همه انسان‌ها عبد خدا \\
\hline

\rowcolor{white}
نگاه به بردگی & طبیعی و ضروری (ارسطو) & مجاز اما مستحب‌الزوال \\
\hline

\rowcolor{EraAncient!15}
منبع آزادی & شهروندی سیاسی & عبودیت خدا (نه انسان دیگر) \\
\hline

\rowcolor{white}
کرامت عمومی & خیر—فقط شهروندان & بله—همه انسان‌ها مکرّم \\
\hline

\rowcolor{EraAncient!15}
مسیر رهایی & بازخرید، آزادسازی & تشویق شرعی به آزادسازی \\
\hline

\rowcolor{white}
آزادی سیاسی & اساسی و هدف & وسیله برای تکلیف الهی \\
\hline

\rowcolor{EraAncient!15}
تناقض اصلی & آزادی بر بردگی بنا شده & سازگاری بردگی با کرامت انسانی \\
\hline

\end{tabularx}
\end{table}

% ─────────────────────────────────────────────────────────────
\subsection{آ.۲ — آزادی و استبداد در حکمرانی: تنش اصلی سنت اسلامی–ایرانی}
\label{subsec:islamtyranny}
% ─────────────────────────────────────────────────────────────

اگر اسلام از یک سو کرامت عمومی را اعلام کرد، از سوی دیگر با معضل پیچیده‌ای روبرو شد: چه کسی و با چه مشروعیتی باید حکومت کند؟

\subsubsection{ساختار کلاسیک: امامت، خلافت، و سلطنت}

سه نهاد تاریخی اصلی در دنیای اسلامی برای حکمرانی وجود داشته:

\begin{enumerate}
  \item \textbf{امامت} (دیدگاه شیعی): امام منصوب از سوی خدا/پیامبر؛ اطاعت از امام واجب است اما امام عادل باید باشد.
  \item \textbf{خلافت} (دیدگاه سنی کلاسیک): خلیفه نماینده‌ی امت است؛ باید عدالت را اجرا کند اما مشروعیتش دست‌کم تا حدی از بیعت می‌آید.
  \item \textbf{سلطنت} (\lr{Mulk}): قدرت واقعی سیاسی که متکلمان اغلب آن را از لحاظ اخلاقی مشکوک می‌دانستند.
\end{enumerate}

\thinker{ابن‌خلدون} در \textit{مقدمه} تحلیل جامعه‌شناختی منحصربه‌فردی از قدرت ارائه داد: قدرت طبیعتاً رو به فساد می‌رود (\lr{asabiyya} رو به زوال است)، و آزادی در نظام سلطانی ذاتاً شکننده است.\footnote{\lr{Ibn Khaldun, \textit{Muqaddimah}, tr. F. Rosenthal (Princeton, 1958), I, pp. 280--295.}} % FIXED: \lr in footnote

\subsubsection{بحث آزادی در اندیشه‌ی متکلمان}

\begin{tcolorbox}[defbox, title={فقه‌الحریه: مفهوم آزادی در فقه و کلام اسلامی}]

\textbf{حرّیت} (آزادی) در فقه اسلامی اول و آخر یک وضعیت حقوقی بود: در برابر عبودیت (بردگی). اما متکلمان آن را به دو حوزه‌ی مهم گسترش دادند:

\textbf{الف. اختیار و جبر:} بحث اصلی کلامی درباره‌ی آزادی اراده. معتزله اختیار انسانی را کامل می‌دانستند؛ اشاعره اراده‌ی الهی را مقدم می‌گرفتند؛ ماتریدیه موضعی میانه داشتند.

\textbf{ب. آزادی سیاسی:} امر به معروف و نهی از منکر به‌عنوان یک نهاد آزادی‌بخش—مشروط که قدرت حاکم را هم در بر می‌گیرد. اما در عمل این نهاد غالباً به حاشیه رانده شد.
\end{tcolorbox}

% ─────────────────────────────────────────────────────────────
\subsection{آ.۳ — بحث ایرانشهری: آزادی در سنت اندیشه‌ی سیاسی ایرانی}
\label{subsec:iranshahri}
% ─────────────────────────────────────────────────────────────

\thinker{داریوش شایگان} و به‌ویژه \thinker{سیدجواد طباطبایی} بحث «ایرانشهری» را مطرح کرده‌اند: آیا در سنت اندیشه‌ی ایرانی، مفهوم آزادی سیاسی به معنای مدرن آن وجود داشته؟

\subsubsection{سنت اندرزنامه‌ای و عدالت پادشاهی}

در سنت اندرزنامه‌نویسی ایرانی (از \textit{خدای‌نامه} تا \textit{سیاست‌نامه‌ی نظام‌الملک})، آزادی مستقیماً مطرح نیست. آنچه هست مفهوم \concept{عدالت} (\lr{dad}) پادشاه است: پادشاه عادل نه از مردم باج ناروا می‌گیرد و نه بی‌دلیل زجر می‌دهد. این نوعی «آزادی از ظلم» است—نه آزادی مشارکت سیاسی.

\begin{tcolorbox}[quotebox]
«پادشاهی و دین دو برادرند که هیچ‌گاه یکدیگر را رها نمی‌کنند. هر زمان که در دین خلل افتاد، در پادشاهی نیز خلل می‌افتد.» — نظام‌الملک، سیاست‌نامه
\end{tcolorbox}

این سنت هرچند استبداد را محدود می‌کند، اما آزادی را از پایین نمی‌سازد—از بالا آن را «اعطا» می‌کند.

\subsubsection{انتقاد طباطبایی: افول اندیشه‌ی سیاسی ایرانی}

\thinker{سیدجواد طباطبایی} استدلال می‌کند که در دوران ورود اسلام، «ایرانشهر»—سنت اندیشه‌ی سیاسی ایرانی که بر پادشاهی آرمانی و عدالت بنا شده بود—با اسلام تلفیق شد، اما این تلفیق ناقص بود: اسلام سیاسی نه از سنت مشارکت یونانی، نه از سنت عدالت ایرانی به‌طور کامل بهره برد.\footnote{طباطبایی، س.ج. (۱۳۸۵). \textit{ابن‌خلدون و علوم اجتماعی}. تهران: ققنوس.}

\subsubsection{طبقه‌بندی مواضع درباره‌ی اسلام و آزادی}

\begin{figure}[h!]
\centering
\caption{طیف‌بندی مواضع اندیشمندان مسلمان و ایرانی درباره‌ی آزادی}
\label{fig:islamspectrum}

\begin{tikzpicture}[xscale=1.1]

% خط اصلی
\draw[line width=2pt, PrimaryDeep, -Stealth]
  (-0.5, 0) -- (11, 0) node[right, font=\small] {\rl{میزان سازگاری اسلام با آزادی سیاسی مدرن}}; % FIXED: \rl

% نقاط روی طیف
\foreach \x/\name/\pos in {
  0.5/{\rl{فقه سنتی}\\% FIXED: گره چندخطی
  \rl{(انتظار حکم شرع)}%
  }/above,
  2.5/{\rl{ابن‌تیمیه}\\%
  \rl{(شریعت=آزادی)}%
  }/below,
  4.5/{\rl{غزالی}\\%
  \rl{(اعتدال کلامی)}%
  }/above,
  6.5/{\rl{اقبال لاهوری}\\%
  \rl{(اجتهاد سیاسی)}%
  }/below,
  8.5/{\rl{عبدالکریم سروش}\\%
  \rl{(قرائت حقوق بشری)}%
  }/above,
  10.0/{\rl{علی شریعتی}\\%
  \rl{(اسلام رهایی‌بخش)}%
  }/below%
}{
  \fill[PrimaryMid] (\x, 0) circle (4pt);
  \node[font=\footnotesize, text width=2.5cm, text centered, \pos]
    at (\x, 0) {\name};
}

% برچسب‌های دو انتها
\node[font=\small\bfseries, color=AccentRed, below] at (0, -0.7) {\rl{سنت‌گرا}}; % FIXED: \rl
\node[font=\small\bfseries, color=AccentGreen, below] at (10.5, -0.7) {\rl{نوگرا/اصلاح‌طلب}}; % FIXED: \rl

% نوار رنگی
\fill[AccentRed!30, opacity=0.4] (0, -0.15) rectangle (3, 0.15);
\fill[AccentGold!30, opacity=0.4] (3, -0.15) rectangle (7, 0.15);
\fill[AccentGreen!30, opacity=0.4] (7, -0.15) rectangle (10.5, 0.15);

\end{tikzpicture}
\end{figure}

\begin{table}[h!]
\centering
\caption{طبقه‌بندی رویکردهای اسلامی به آزادی سیاسی}
\label{tab:islamfreedom}
\renewcommand{\arraystretch}{1.5}

\begin{tabularx}{\textwidth}{|>{\bfseries}p{2.5cm}|>{\raggedright}X|>{\raggedright}X|>{\raggedright\arraybackslash}X|}
\hline
\rowcolor{PrimaryDeep}
\textcolor{white}{\textbf{رویکرد}} &
\textcolor{white}{\textbf{تعریف آزادی}} &
\textcolor{white}{\textbf{نمایندگان}} &
\textcolor{white}{\textbf{نقد اصلی}} \\
\hline

\rowcolor{EraAncient!15}
فقه سنتی & آزادی=اجرای احکام شرع & فقهای کلاسیک & ناتوانی در پاسخ به دولت مدرن \\
\hline

\rowcolor{EraMedieval!15}
اسلام سیاسی & آزادی=اجرای شریعت توسط دولت اسلامی & قطب، مودودی & خلط دین و سیاست \\
\hline

\rowcolor{EraEarlyMod!15}
اصلاح‌طلبی کلاسیک & آزادی از استبداد شرقی، پذیرش نهادهای مدرن & عبده، کواکبی & سازگاری ناقص \\
\hline

\rowcolor{EraModern!15}
رهایی‌بخشی اسلامی & آزادی=رهایی از استعمار و استثمار & شریعتی، قرضاوی & آمیختگی با ناسیونالیسم \\
\hline

\rowcolor{EraContemp!15}
نوهرمنوتیک & آزادی=بازخوانی نصوص با تأکید بر کرامت & سروش، شبستری & متهم به سکولاریسم \\
\hline

\rowcolor{EraPostmod!15}
اسلام لیبرال & جدایی دین از دولت، حقوق بشر جهانی & المسیری، النعیم & کم‌پایگاه در جوامع مسلمان \\
\hline

\end{tabularx}
\end{table}

\newpage
